\documentclass[11pt,a4paper]{article}
\usepackage[T1]{fontenc}
\usepackage[utf8]{inputenc}
\usepackage{lmodern}
\usepackage[a4paper,margin=25mm,headheight=15pt]{geometry}
\usepackage{amsmath,amssymb,amsthm,mathtools}
\usepackage{microtype,booktabs,longtable,array,xcolor,listings,fancyhdr}
\usepackage[hyphens]{url}
\usepackage{xurl}
\usepackage[colorlinks=true,linkcolor=blue!45!black,citecolor=blue!45!black,urlcolor=blue!45!black]{hyperref}
\hypersetup{pdftitle={A smaller iPRA presentation by first-instance witness discharge},pdfauthor={Research report prepared for Vladimir Reshetnikov}}
\newtheorem{theorem}{Theorem}[section]
\newtheorem{lemma}[theorem]{Lemma}
\newtheorem{proposition}[theorem]{Proposition}
\newtheorem{corollary}[theorem]{Corollary}
\theoremstyle{definition}\newtheorem{definition}[theorem]{Definition}
\theoremstyle{remark}\newtheorem{remark}[theorem]{Remark}
\newcommand{\ipra}{\mathsf{iPRA}}
\newcommand{\nf}{\operatorname{nf}}
\newcommand{\FV}{\operatorname{FV}}
\pagestyle{fancy}\fancyhf{}
\fancyhead[L]{\small First-instance witness discharge for iPRA}
\fancyhead[R]{\small\thepage}
\renewcommand{\headrulewidth}{0.3pt}
\setlength{\emergencystretch}{2em}
\allowdisplaybreaks[2]
\raggedbottom
\lstset{language=Python,basicstyle=\footnotesize\ttfamily,columns=fullflexible,breaklines=true,breakatwhitespace=false,showstringspaces=false,frame=single,rulecolor=\color{black!20},tabsize=4}
\title{A smaller presentation of intuitionistic\newline primitive recursive arithmetic\\[4pt]
\large First-instance witness discharge}
\author{Research report prepared for Vladimir Reshetnikov}
\date{11 November 2026}
\begin{document}
\maketitle
\begin{abstract}
The current witness-discharge presentation of full first-order iPRA has
three inference families. Its elimination rule takes a major formula,
a local equality proof $b=t$, and a local formula proof $F$, then exports
the normalized instance $F[t/b]$. We prove that the equality proof is
unnecessary as a primitive premise: a deterministic test of the opened
clause's zero and first instances supplies the only substitution that
the rule needs to perform. The new rule has two proof premises and no
replacement-term argument. Nevertheless it preserves every open sequent,
not merely the closed equations.

The reconstruction obtains ordinary elimination and numerical cases
before equality, and then recovers arbitrary transport through a
two-row naming matrix. The inherited full-iPRA bi-interpretation therefore
has both numerical domains whole, actual equality, and identity numerical
composites. We include the complete primitive specification, proofs,
the normalization and arithmetic bridges, a standalone Python checker,
proof translations, adversarial examples, and reproducible measurements.
The result is a concrete simplification relative to the current
witness-discharge presentation, not an absolute minimality theorem.
\end{abstract}
\tableofcontents
\clearpage

\section{Result and comparison boundary}
\label{sec:boundary}

The comparison repository is ProveIt~\cite{baseline}, under the subtree
\begin{center}\small\path{Logic/PrimitiveRecursiveArithmetic/ConciseIPRA}.\end{center}
The pinned commit is
\begin{center}\small\ttfamily 17898ea7e6a1300dfdcad71040d5daeed67916d4.\end{center}
Its README lists sixty-seven filed research packages. We also inspected
the six archives in \path{docs/incoming}, including the recent
\texttt{witness\_discharge} and \texttt{uniform\_binding} packages.
The immediate mathematical control is the incoming
\path{ipra_witness_discharge_2026-11-11.zip}~\cite{witness}.
In particular, three inference families and an initial reflexivity clause
are already present there. The new theorem concerns deletion of one
\emph{proof premise} and deterministic recovery of its substitution.

The complete proposed calculus is denoted $\mathcal A$.
Its formulas have only equality and the paired binder $D$; its numerical
terms are canonical scalar iterators. There is one initial clause and
three inference families: Use, Intro, and Name. The standard reading is
\[
  D_{x,y}B(x,y)
    \quad\rightsquigarrow\quad
  \exists x\bigl(B(x,0)\to\forall y\,B(x,y)\bigr).
\]
The syntactic reductions used to obtain canonical terms are explicitly
part of the theory's specification and of the measured checker.

\paragraph{Reference theory.}
Here $\ipra$ means full intuitionistic first-order logic with equality,
zero, successor, and a function symbol for every fixed primitive-recursive
description, the usual defining equations and successor axioms, and
induction restricted to quantifier-free formulas, with parameters.
It is not the quantifier-free proof language alone. The equivalent
atomic-induction convention in the full primitive-recursive language is
also used in the literature~\cite{khaniki}; the equality-reflection
appendix gives the required equivalence in this development.

\paragraph{What is proved.}
For every canonical open clause sequent, the old witness-discharge
calculus $\mathcal W$ and $\mathcal A$ prove exactly the same conclusion
from exactly the same assumptions. The translation is the identity on
their numerical and formula languages. Composing with the scalar
arithmetic bridge proves full bi-interpretability and equiconsistency
with $\ipra$. The two numerical composite interpretations are internally
isomorphic through the graph $x=y$; both formula composites are also
proved, so this is stronger than a standard-model agreement or an
equational conservativity statement.

\paragraph{What is inherited.}
The scalar iterator, the $D$ binder, the first-order logical macros,
full computation normalization, the two-row naming method, predecessor,
equality reflection, and the primitive-recursive compiler are inherited
from the repository. The new contribution is the completeness of a
fixed first-instance restriction on witness discharge. The transparent
certificate representation and its complete implementation are also
given, with matched controls. No universal priority or absolute
minimum-source-size claim is made.

\paragraph{How to read the report.}
Section~\ref{sec:deterministic-name} contains the complete native
axiomatization and the new proof. Section~\ref{sec:canonical-bridge}
connects canonical syntax to the raw scalar basis. The implementation
section records the complete checker contract and its measured scope.
The appendices reproduce the attributed logical and arithmetic
foundations, so the full-iPRA conclusion does not rest on an omitted
predecessor or pairing argument.
